% A DELIMITER SCAN STEPS OVER \relax; A NUMBER SCAN DOES NOT.
%
% The reference asks for "the next non-blank non-relax non-call token"
% where it wants a delimiter -- after \left, after \right, and for both
% delimiters of \atopwithdelims and its kin. So \relax standing there is
% skipped, and what is reported when no delimiter follows is the token
% AFTER the \relax:
%
%   $\left\relax\mathchardef ...   <to be read again> \mathchardef
%
% A number scan does not step over \relax: \relax ends it, and the
% "Missing number" that follows names the \relax itself. \radical and
% \mathaccent read numbers, so
%
%   $\radical\relax"270370 x$      <to be read again> \relax
%
% and the \relax put back there is skipped by the SUBFORMULA scan that
% comes next, which asks for a non-blank non-relax non-call token too --
% so it is never obeyed and never traced.
%
% A delimiter that is not one at all -- INITEX gives every character a
% \delcode of -1 but for `.' -- is put back after the report, so the
% character is read again and traced as itself.
%
% ALSO: a definition given something that cannot be a name puts that
% token back and INSERTS \inaccessible in front of it. The context names
% the inserted one, and the definition then reads it as it would read any
% other name.
\catcode`\{=1 \catcode`\}=2 \catcode`\$=3 \catcode`\^=7 \catcode`\_=8
\tracingonline=1 \tracingcommands=2
\message{[A \left and \right with \relax before the delimiter]}
$\left\relax(x\right\relax)$
\message{[B a non-delimiter after \relax]}
$\left\relax\mathchardef x\right.$
\message{[C \atopwithdelims]}
$1\atopwithdelims\relax(\relax)2$
\message{[D \radical reads a number, which \relax ends]}
$\radical\relax"270370 x$
\message{[E \mathaccent, the same]}
$\mathaccent\relax"7016 x$
\message{[F \delimiter, the same]}
$\delimiter\relax"0 x$
\message{[done]}
\end
